% Part: first-order-logic
% Chapter: beyond
% Section: introduction

\documentclass[../../../include/open-logic-section]{subfiles}

\begin{document}

\olfileid[fr]{fol}{byd}{int}

\olsection{Vue d'ensemble}

La logique du premier ordre n'est pas le seul système logique qui
présente un intérêt : elle connaît de nombreuses extensions et
variantes. Une logique comprend généralement la spécification
formelle d'un langage, le plus souvent, mais pas toujours, un système
déductif et, le plus souvent aussi, une sémantique visée. Cet emploi
technique du mot soulève cependant une question : quel rapport les
logiques autres que celle du premier ordre entretiennent-elles avec
la « logique » au sens intuitif ou philosophique ? Tous les systèmes
décrits ci-dessous visent à modéliser une forme ou une autre de
raisonnement. Peut-on dire ce qui en fait des systèmes logiques ?

La réponse n'est pas simple. Le mot «~logique~» a des usages différents
selon le contexte. La notion de logique, comme celle de vérité, a fait
l'objet d'analyses philosophiques très diverses. On pourrait, par exemple, assigner au raisonnement
logique la tâche de déterminer quels énoncés sont nécessairement
vrais, vrais \emph{a priori}, vrais indépendamment de
l'interprétation des termes non logiques, vrais en vertu de leur
forme ou vrais par convention linguistique. Chacune de ces
conceptions exige pourtant de nombreuses précisions. Même si l'on
s'en tient à la logique employée en mathématiques, sa portée ne fait
guère l'unanimité. Ainsi, dans les \textit{Principia Mathematica},
Russell et Whitehead ont tenté de développer les mathématiques sur la
base de la logique, dans la tradition {\em logiciste} inaugurée par
Frege. Leur système était une forme de logique d'ordre supérieur,
semblable à celle qui sera décrite plus loin. Ils ont finalement dû
introduire des axiomes qui, selon la plupart des critères, ne semblent
pas purement logiques, notamment l'axiome de l'infini et l'axiome de
réductibilité. On peut néanmoins soutenir que certaines formes de
raisonnement d'ordre supérieur doivent être reconnues comme logiques.
À l'inverse, Quine, dont l'ontologie n'admet pas les « propositions »
comme objets légitimes de discours, affirme que les logiques du
second ordre et d'ordre supérieur sont, en réalité, des formes de
théorie des ensembles déguisées. Autrement dit, les systèmes qui
quantifient sur les prédicats ne seraient pas purement logiques.

Laissons pour l'instant ces questions philosophiques de côté et
considérons simplement les systèmes ci-dessous comme des
idéalisations formelles de diverses formes de raisonnement, qu'on
les juge logiques ou non.

\end{document}
